\par\noindent\textbf{(1)} 小斜率下弯曲储能为(EI/2)∫w''²dx，压缩载荷释放功为(P/2)∫w'²dx，故二阶总势Π2=(1/2)∫[EIw''²−Pw'²]dx。变分并用铰支边界积分分部得EI w''''+P w''=0；Euler–Lagrange方程是空间四阶，不能直接写成简谐振子的时间方程。
\par\noindent\textbf{(2)} 取本征挠度w=A sin(nπx/L)，满足位移零和弯矩EIw''=0的两端条件。代入得P\_n=EI(nπ/L)²，n=1,2,…。最先失稳为P\_cr=π²EI/L²，说明加大弯曲刚度或缩短杆长会大幅提高承载阈值。
\par\noindent\textbf{(3)} 将一般w按正弦完备基展开，Π2=(L/4)Σ\_n A\_n²[EI(nπ/L)\textasciicircum{}4−P(nπ/L)²]。P<P\_cr时每项正，直杆局部稳定；P>P\_cr时n=1方向二次型负，任意小首模挠度即可降能。P=P\_cr二次项为零，需要高阶几何能量确定分岔。
\par\noindent\textbf{(4)} 若固定端缩短量控制，可保留中心线伸长引起的正四次项C A⁴（C>0），首模约化能量成为Π≈c2(P\_cr−P)A²+C A⁴。P>P\_cr时稳态幅值A∝√(P−P\_cr)，但系数依加载方式与轴向刚度，不能由只含二次项的线性屈曲方程给出。
